\documentclass[11pt, letterpaper]{article}

%% ── Packages ──────────────────────────────────────────────────────────────
\usepackage[utf8]{inputenc}
\usepackage[T1]{fontenc}
\usepackage{lmodern}
\usepackage[margin=1in]{geometry}
\usepackage{graphicx}
\usepackage{xcolor}
\usepackage{listings}
\usepackage{booktabs}
\usepackage{longtable}
\usepackage{hyperref}
\usepackage{amsmath, amssymb}
\usepackage{enumitem}
\usepackage{fancyhdr}
\usepackage{titlesec}
\usepackage{caption}
\usepackage{tcolorbox}
\usepackage{array}
\usepackage{tabularx}

%% ── Colour palette ───────────────────────────────────────────────────────
\definecolor{codegreen}{rgb}{0.25,0.49,0.24}
\definecolor{codegray}{rgb}{0.5,0.5,0.5}
\definecolor{codepurple}{rgb}{0.49,0.0,0.65}
\definecolor{codebg}{rgb}{0.97,0.97,0.97}
\definecolor{highsev}{HTML}{D32F2F}
\definecolor{medsev}{HTML}{F57C00}
\definecolor{lowsev}{HTML}{388E3C}
\definecolor{accentblue}{HTML}{1565C0}

%% ── Code listing style ──────────────────────────────────────────────────
\lstdefinestyle{pythonstyle}{
    backgroundcolor=\color{codebg},
    commentstyle=\color{codegreen}\itshape,
    keywordstyle=\color{accentblue}\bfseries,
    stringstyle=\color{codepurple},
    basicstyle=\ttfamily\footnotesize,
    breaklines=true,
    breakatwhitespace=false,
    numbers=left,
    numberstyle=\tiny\color{codegray},
    numbersep=8pt,
    tabsize=4,
    frame=single,
    rulecolor=\color{codegray!40},
    showstringspaces=false,
    captionpos=b,
    xleftmargin=1.5em,
    framexleftmargin=1.5em,
}
\lstset{style=pythonstyle, language=Python}

%% ── YAML listing style ─────────────────────────────────────────────────
\lstdefinestyle{yamlstyle}{
    backgroundcolor=\color{codebg},
    basicstyle=\ttfamily\footnotesize,
    breaklines=true,
    numbers=left,
    numberstyle=\tiny\color{codegray},
    numbersep=8pt,
    frame=single,
    rulecolor=\color{codegray!40},
    showstringspaces=false,
    xleftmargin=1.5em,
    framexleftmargin=1.5em,
    commentstyle=\color{codegreen}\itshape,
    morecomment=[l]{\#},
}

%% ── Severity boxes ──────────────────────────────────────────────────────
\newtcolorbox{highbox}{colback=highsev!8, colframe=highsev!80, title=\textbf{HIGH Severity}, fonttitle=\bfseries\color{white}, coltitle=white, colbacktitle=highsev}
\newtcolorbox{medbox}{colback=medsev!8, colframe=medsev!80, title=\textbf{MEDIUM Severity}, fonttitle=\bfseries\color{white}, coltitle=white, colbacktitle=medsev}
\newtcolorbox{lowbox}{colback=lowsev!8, colframe=lowsev!80, title=\textbf{LOW Severity}, fonttitle=\bfseries\color{white}, coltitle=white, colbacktitle=lowsev}

%% ── Header / footer ────────────────────────────────────────────────────
\pagestyle{fancy}
\fancyhf{}
\lhead{\small\textsc{Price Impact HTE Pipeline}}
\rhead{\small Code Report}
\cfoot{\thepage}
\renewcommand{\headrulewidth}{0.4pt}

%% ── Section formatting ─────────────────────────────────────────────────
\titleformat{\section}{\Large\bfseries\color{accentblue}}{\thesection}{1em}{}
\titleformat{\subsection}{\large\bfseries\color{accentblue!80}}{\thesubsection}{1em}{}
\titleformat{\subsubsection}{\normalsize\bfseries}{\thesubsubsection}{1em}{}

\hypersetup{
    colorlinks=true,
    linkcolor=accentblue,
    urlcolor=accentblue!80,
    citecolor=accentblue,
}

%% ═════════════════════════════════════════════════════════════════════════
\begin{document}

%% ── Title page ──────────────────────────────────────────────────────────
\begin{titlepage}
    \centering
    \vspace*{2cm}
    {\Huge\bfseries Price Impact HTE\\Analysis Pipeline\par}
    \vspace{0.8cm}
    {\LARGE Code Review \& Error Analysis Report\par}
    \vspace{2cm}
    {\Large Brian Butler\par}
    \vspace{0.5cm}
    {\large June 3, 2026\par}
    \vfill
    {\large
    \begin{tabular}{rl}
        \textbf{Project:}  & \texttt{Market/} \\
        \textbf{Language:} & Python 3.13 \\
        \textbf{Libraries:}& econml, scikit-learn, DoWhy, pandas, NumPy \\
        \textbf{Files:}    & 8 source files, 718 total lines \\
    \end{tabular}
    }
    \vspace{1cm}
\end{titlepage}

%% ── Table of Contents ───────────────────────────────────────────────────
\tableofcontents
\newpage

%% ═════════════════════════════════════════════════════════════════════════
\section{Project Overview}
%% ═════════════════════════════════════════════════════════════════════════

This project implements a \textbf{Heterogeneous Treatment Effect (HTE)} analysis pipeline for measuring \textbf{price impact} in financial markets. It:

\begin{enumerate}[leftmargin=2em]
    \item Generates synthetic trade data with a known ground-truth CATE,
    \item Engineers liquidity and volatility features,
    \item Fits five causal ML estimators (CausalForest + four meta-learners),
    \item Evaluates estimator accuracy against the known ground truth, and
    \item Produces publication-quality visualizations.
\end{enumerate}

The ground-truth Conditional Average Treatment Effect is defined analytically as:
\begin{equation}
    \tau(\mathbf{X}) \;=\; \beta_0 \;+\; \frac{\beta_1}{\text{depth}} \;+\; \beta_2 \cdot \text{volatility}
    \label{eq:cate}
\end{equation}

\noindent where $\beta_0 = 0.0005$, $\beta_1 = 0.5$, and $\beta_2 = 0.05$.

\subsection{Data Flow}

The pipeline executes six sequential steps, orchestrated by \texttt{main.py}:

\begin{center}
\begin{tabular}{clll}
    \toprule
    \textbf{Step} & \textbf{Module} & \textbf{Action} & \textbf{Output} \\
    \midrule
    1 & \texttt{dgp.py}           & Generate synthetic data      & DataFrame (200k rows) \\
    2 & \texttt{features.py}      & Engineer features             & Standardised covariates \\
    3 & \texttt{estimators.py}    & Fit 5 HTE estimators          & Trained model objects \\
    4 & \texttt{estimators.py}    & Predict CATE on test set      & CATE arrays + CIs \\
    5 & \texttt{evaluation.py}    & Compute accuracy metrics      & Results DataFrame \\
    6 & \texttt{visualization.py} & Generate plots                & PNG files in \texttt{results/} \\
    \bottomrule
\end{tabular}
\end{center}

\subsection{File Index}

\begin{center}
\begin{tabular}{llrl}
    \toprule
    \textbf{File} & \textbf{Path} & \textbf{Lines} & \textbf{Role} \\
    \midrule
    \texttt{params.yaml}       & \texttt{config/}  & 66  & All hyperparameters \\
    \texttt{dgp.py}            & \texttt{src/}     & 98  & Synthetic data generation \\
    \texttt{features.py}       & \texttt{src/}     & 86  & Feature engineering \\
    \texttt{causal\_graph.py}  & \texttt{src/}     & 69  & Causal DAG \& assumptions \\
    \texttt{estimators.py}     & \texttt{src/}     & 129 & Five HTE estimators \\
    \texttt{evaluation.py}     & \texttt{src/}     & 127 & Metrics \& GATES analysis \\
    \texttt{visualization.py}  & \texttt{src/}     & 193 & All plotting functions \\
    \texttt{main.py}           & root              & 85  & Pipeline orchestration \\
    \bottomrule
\end{tabular}
\end{center}


%% ═════════════════════════════════════════════════════════════════════════
\newpage
\section{Configuration --- \texttt{config/params.yaml}}
\label{sec:config}
%% ═════════════════════════════════════════════════════════════════════════

All hyperparameters are centralised in a single YAML file.
The configuration is loaded by \texttt{dgp.load\_config()} and threaded through every module.

\begin{lstlisting}[style=yamlstyle, caption={Full configuration file.}]
random_seed: 42

dgp:
  n_observations: 200000
  train_fraction: 0.7

  # Covariate distributions
  spread_loc: -2.0
  spread_scale: 0.6
  depth_scale: 50.0
  volatility_scale: 0.02

  # Treatment (confounded by market state)
  treatment_noise_scale: 5.0
  treatment_depth_coef: 0.3
  treatment_vol_coef: -10.0

  # Outcome -- true CATE: tau(X) = B0 + B1/depth + B2*vol
  baseline_drift_scale: 0.001
  cate_intercept: 0.0005
  cate_depth_coef: 0.5
  cate_vol_coef: 0.05
  outcome_noise_scale: 0.002

  # AR(1) autocorrelation in noise (0 = IID)
  ar1_coefficient: 0.3

features:
  volatility_window: 20
  winsorize_percentile: 0.01

estimators:
  causal_forest:
    n_estimators: 1000
    min_samples_leaf: 20
    max_depth: null
    discrete_treatment: false

  s_learner:
    base_model: "GradientBoostingRegressor"
    n_estimators: 200

  t_learner:
    base_model: "GradientBoostingRegressor"
    n_estimators: 200

  x_learner:
    base_model: "GradientBoostingRegressor"
    n_estimators: 200

  dr_learner:
    base_model: "GradientBoostingRegressor"
    n_estimators: 200

  binary_treatment_quantile: 0.75

evaluation:
  n_bootstrap: 200
  ci_level: 0.90
  n_cate_bins: 5

output:
  results_dir: "results"
  figure_format: "png"
  figure_dpi: 300
\end{lstlisting}

\subsection{Error Analysis}

\begin{center}
\begin{tabularx}{\textwidth}{c c c X}
    \toprule
    \textbf{ID} & \textbf{Param} & \textbf{Sev.} & \textbf{Description} \\
    \midrule
    C1 & \texttt{n\_bootstrap} & \textcolor{lowsev}{Low} & Declared but never referenced in any code. Dead config. \\
    C2 & \texttt{figure\_dpi}  & \textcolor{lowsev}{Low} & Never read; \texttt{visualization.py} hardcodes DPI to 300. \\
    C3 & \texttt{load\_config} & \textcolor{medsev}{Med} & No schema validation. A config typo produces a cryptic \texttt{KeyError}. \\
    \bottomrule
\end{tabularx}
\end{center}


%% ═════════════════════════════════════════════════════════════════════════
\newpage
\section{Data Generation --- \texttt{src/dgp.py}}
\label{sec:dgp}
%% ═════════════════════════════════════════════════════════════════════════

\subsection{Purpose}

Generate synthetic trade data where the ground-truth CATE is known analytically (Equation~\ref{eq:cate}), allowing precise evaluation of estimator accuracy.

\subsection{Key Functions}

\subsubsection{\texttt{load\_config()}}
Reads the YAML configuration file and returns a nested dictionary.

\begin{lstlisting}[caption={Configuration loader.}]
def load_config(config_path: str = "config/params.yaml") -> dict:
    """Load the YAML configuration file."""
    with open(config_path, "r") as f:
        return yaml.safe_load(f)
\end{lstlisting}

\subsubsection{\texttt{compute\_true\_cate()}}
Implements Equation~\ref{eq:cate} with a small epsilon guard against division by zero.

\begin{lstlisting}[caption={Ground-truth CATE computation.}]
def compute_true_cate(depth, volatility, cfg) -> np.ndarray:
    dgp = cfg["dgp"]
    cate = (
        dgp["cate_intercept"]
        + dgp["cate_depth_coef"] / (depth + 1e-6)
        + dgp["cate_vol_coef"] * volatility
    )
    return cate
\end{lstlisting}

\subsubsection{\texttt{generate\_data()}}
Creates $n = 200{,}000$ observations with:
\begin{itemize}[leftmargin=2em]
    \item \textbf{Covariates:} spread $\sim \text{LogNormal}$, depth $\sim \text{Exp}(50)$, volatility $\sim |\mathcal{N}(0, 0.02)|$
    \item \textbf{Treatment:} trade\_size $= 0.3 \cdot \text{depth} - 10 \cdot \text{volatility} + \varepsilon$ (confounded)
    \item \textbf{Outcome:} $Y = g(\mathbf{X}) + \tau(\mathbf{X}) \cdot T + \varepsilon_{\text{AR}(1)}$
\end{itemize}

\begin{lstlisting}[caption={Core data generation logic (abbreviated).}]
# Market-state covariates
spread     = rng.lognormal(mean=spread_loc, sigma=spread_scale, size=n)
depth      = rng.exponential(scale=depth_scale, size=n)
volatility = np.abs(rng.normal(0, volatility_scale, size=n))

# Treatment -- confounded by market state
trade_size = np.maximum(depth_coef * depth + vol_coef * vol + noise, 1.0)

# Outcome with optional AR(1) autocorrelation in noise
price_change = baseline_drift + true_cate * trade_size + noise
\end{lstlisting}

\subsubsection{\texttt{split\_data()}}
Performs a temporal (positional) train/test split at 70/30.

\begin{lstlisting}[caption={Train/test split.}]
def split_data(df, cfg):
    frac = cfg["dgp"]["train_fraction"]
    split_idx = int(len(df) * frac)
    return df.iloc[:split_idx].copy(), df.iloc[split_idx:].copy()
\end{lstlisting}

\subsection{Error Analysis}

\begin{center}
\begin{tabularx}{\textwidth}{c c c X}
    \toprule
    \textbf{ID} & \textbf{Lines} & \textbf{Sev.} & \textbf{Description} \\
    \midrule
    D1 & 79--83 & \textcolor{medsev}{Med} &
        \texttt{split\_data} is called ``temporal'' but \texttt{generate\_data} produces rows with no time index. The split relies on row order never being shuffled---a fragile assumption. \\
    D2 & 21 & \textcolor{lowsev}{Low} &
        \texttt{depth} drawn from Exp(50) can be near zero. The \texttt{1e-6} guard works, but very small depths produce enormous CATEs that may dominate evaluation metrics. \\
    D3 & (eval 96) & \textcolor{medsev}{Med} &
        \texttt{cfg.copy()} is shallow---nested dicts like \texttt{cfg["dgp"]} still reference the original. If any downstream code mutated \texttt{cfg\_copy["dgp"]}, it would corrupt the original config. Currently safe but latent. \\
    D4 & --- & \textcolor{lowsev}{Low} &
        No \texttt{\_\_init\_\_.py} in \texttt{src/}. The \texttt{sys.path} hack in \texttt{main.py} works around this, but relative imports (\texttt{from .dgp}) used in \texttt{evaluation.py} will fail. \\
    \bottomrule
\end{tabularx}
\end{center}


%% ═════════════════════════════════════════════════════════════════════════
\newpage
\section{Feature Engineering --- \texttt{src/features.py}}
\label{sec:features}
%% ═════════════════════════════════════════════════════════════════════════

\subsection{Purpose}
Transform raw covariates into modelling features: winsorise extremes, derive liquidity ratios, add rolling volatility, and Z-score standardise.

\subsection{Key Functions}

\subsubsection{\texttt{winsorize()}}
\begin{lstlisting}[caption={Winsorisation: clip extreme values at tail percentiles.}]
def winsorize(series, pct=0.01):
    lower = series.quantile(pct)
    upper = series.quantile(1 - pct)
    return series.clip(lower, upper)
\end{lstlisting}

\subsubsection{\texttt{add\_liquidity\_features()}}
\begin{lstlisting}[caption={Derived liquidity proxies.}]
def add_liquidity_features(df):
    df = df.copy()
    df["log_spread"]         = np.log1p(df["spread"])
    df["inv_depth"]          = 1.0 / (df["depth"] + 1e-6)
    df["spread_depth_ratio"] = df["spread"] / (df["depth"] + 1e-6)
    return df
\end{lstlisting}

\subsubsection{\texttt{standardize()}}
Z-score normalisation that can reuse training statistics for the test set.
\begin{lstlisting}[caption={Z-score standardisation with fit/transform pattern.}]
def standardize(df, columns, fit_stats=None):
    df = df.copy()
    if fit_stats is None:
        fit_stats = {}
        for col in columns:
            fit_stats[col] = {"mean": df[col].mean(),
                              "std":  df[col].std()}
    for col in columns:
        mu    = fit_stats[col]["mean"]
        sigma = fit_stats[col]["std"]
        df[col] = (df[col] - mu) / (sigma + 1e-10)
    return df, fit_stats
\end{lstlisting}

\subsubsection{\texttt{engineer\_features()} --- Full Pipeline}
\begin{lstlisting}[caption={Master feature engineering function (abbreviated).}]
def engineer_features(df, cfg, fit_stats=None):
    pct = cfg["features"]["winsorize_percentile"]
    for col in ["spread", "depth", "volatility", "trade_size"]:
        df[col] = winsorize(df[col], pct)        # BUG: mutates input
    df = add_liquidity_features(df)
    df = add_rolling_volatility(df, window=...)
    df, fit_stats = standardize(df, covariate_cols, fit_stats)
    return df, fit_stats
\end{lstlisting}

\subsection{Error Analysis}

\begin{center}
\begin{tabularx}{\textwidth}{c c c X}
    \toprule
    \textbf{ID} & \textbf{Lines} & \textbf{Sev.} & \textbf{Description} \\
    \midrule
    F1 & 59--60 & \textcolor{highsev}{\textbf{High}} &
        \textbf{Mutates input DataFrame.} \texttt{engineer\_features} applies \texttt{winsorize} to columns \emph{before} making a copy. The caller's original DataFrame (\texttt{train\_raw}, \texttt{test\_raw}) is silently modified in-place. \\[6pt]
    F2 & 18--23 & \textcolor{highsev}{\textbf{High}} &
        \textbf{Rolling volatility uses the outcome variable} (\texttt{price\_change}). In a real setting this would be data leakage. In this synthetic setup the feature is computed but not used by the estimators, so it wastes computation and creates confusion. \\[6pt]
    F3 & 65--68 & \textcolor{medsev}{Med} &
        \textbf{Feature mismatch.} Seven columns are standardised, but \texttt{main.py} only feeds three (\texttt{spread}, \texttt{depth}, \texttt{volatility}) to the estimators. The extra four features are never used. \\[6pt]
    F4 & 59--60 & \textcolor{medsev}{Med} &
        \textbf{Test-set quantile leakage.} \texttt{winsorize} recalculates quantile bounds from the test data instead of reusing training bounds. Inconsistent with the \texttt{standardize} function which correctly reuses \texttt{fit\_stats}. \\
    \bottomrule
\end{tabularx}
\end{center}


%% ═════════════════════════════════════════════════════════════════════════
\newpage
\section{Causal Graph --- \texttt{src/causal\_graph.py}}
\label{sec:causalgraph}
%% ═════════════════════════════════════════════════════════════════════════

\subsection{Purpose}
Encodes the assumed causal DAG and the four identifying assumptions. Provides a DoWhy \texttt{CausalModel} builder for refutation tests.

\subsection{Causal DAG}

The assumed causal structure is:

\begin{center}
\begin{tabular}{lcl}
    spread      & $\longrightarrow$ & trade\_size \\
    depth       & $\longrightarrow$ & trade\_size \\
    volatility  & $\longrightarrow$ & trade\_size \\[4pt]
    spread      & $\longrightarrow$ & price\_change \\
    depth       & $\longrightarrow$ & price\_change \\
    volatility  & $\longrightarrow$ & price\_change \\
    trade\_size & $\longrightarrow$ & price\_change \\
\end{tabular}
\end{center}

\noindent All three market-state variables are \textbf{common causes} (confounders) of both treatment and outcome.

\subsection{Identifying Assumptions}

\begin{enumerate}[leftmargin=2em]
    \item \textbf{Unconfoundedness:} Conditional on (spread, depth, volatility), \texttt{trade\_size} is as-good-as-random. No unobserved confounders affect both $T$ and $Y$.
    \item \textbf{Overlap (Positivity):} For every market state $\mathbf{X}$, there is positive probability of observing any trade size $T$.
    \item \textbf{SUTVA:} One trade's outcome does not depend on other trades' treatment assignments.
    \item \textbf{Consistency:} The observed outcome under treatment $T = t$ equals the potential outcome $Y(t)$.
\end{enumerate}

\subsection{Error Analysis}

\begin{center}
\begin{tabularx}{\textwidth}{c c c X}
    \toprule
    \textbf{ID} & \textbf{Lines} & \textbf{Sev.} & \textbf{Description} \\
    \midrule
    G1 & --- & \textcolor{lowsev}{Low} &
        Module is never imported in \texttt{main.py}. DoWhy refutation tests are defined but never run. Assumptions are documented but not programmatically verified. \\
    G2 & 48--50 & \textcolor{lowsev}{Low} &
        \texttt{build\_dowhy\_model} silently returns \texttt{None} if \texttt{dowhy} is not installed. Any caller that doesn't check for \texttt{None} would crash. \\
    \bottomrule
\end{tabularx}
\end{center}


%% ═════════════════════════════════════════════════════════════════════════
\newpage
\section{Estimators --- \texttt{src/estimators.py}}
\label{sec:estimators}
%% ═════════════════════════════════════════════════════════════════════════

\subsection{Purpose}
Wrap five HTE estimators from the \texttt{econml} library in a unified \texttt{EstimatorSuite} class.

\subsection{Treatment Binarisation}

The meta-learners require a binary treatment. Continuous \texttt{trade\_size} is converted at the 75th percentile:

\begin{lstlisting}[caption={Binarise continuous treatment at a quantile threshold.}]
def binarize_treatment(trade_size, quantile=0.75):
    threshold = np.quantile(trade_size, quantile)
    return (trade_size >= threshold).astype(int)
\end{lstlisting}

\begin{medbox}
CausalForest uses the \textbf{continuous} treatment ($\partial Y / \partial T$), while all four meta-learners use a \textbf{binarised} treatment ($\mathbb{E}[Y(1) - Y(0)]$). These are \emph{different causal quantities}. Comparing their RMSE against the same ground-truth $\tau(\mathbf{X})$ is scientifically problematic (see issue E2).
\end{medbox}

\subsection{EstimatorSuite Class}

\begin{center}
\begin{tabularx}{\textwidth}{l X}
    \toprule
    \textbf{Method} & \textbf{Description} \\
    \midrule
    \texttt{fit\_all()}            & Fits all five estimators on training data \\
    \texttt{\_fit\_causal\_forest()} & \texttt{econml.dml.CausalForestDML} --- continuous treatment \\
    \texttt{\_fit\_s\_learner()}     & \texttt{econml.metalearners.SLearner} --- single model \\
    \texttt{\_fit\_t\_learner()}     & \texttt{econml.metalearners.TLearner} --- separate models per arm \\
    \texttt{\_fit\_x\_learner()}     & \texttt{econml.metalearners.XLearner} --- cross-fitted imputation \\
    \texttt{\_fit\_dr\_learner()}    & \texttt{econml.metalearners.DRLearner} --- doubly robust \\
    \texttt{estimate\_cate()}      & Predicts CATE on new covariates for all estimators \\
    \texttt{confidence\_intervals()} & Returns CIs where supported (e.g.\ CausalForest) \\
    \bottomrule
\end{tabularx}
\end{center}

\subsection{Error Analysis}

\begin{center}
\begin{tabularx}{\textwidth}{c c c X}
    \toprule
    \textbf{ID} & \textbf{Lines} & \textbf{Sev.} & \textbf{Description} \\
    \midrule
    E1 & 102 & \textcolor{highsev}{\textbf{High}} &
        \textbf{Wrong import path.} \texttt{from econml.metalearners import DRLearner} is likely incorrect. In standard \texttt{econml}, \texttt{DRLearner} lives in \texttt{econml.dr}, not \texttt{econml.metalearners}. This will raise \texttt{ImportError} at runtime. \\[6pt]
    E2 & 40--49 & \textcolor{highsev}{\textbf{High}} &
        \textbf{Comparing apples to oranges.} CausalForest estimates a marginal effect of continuous treatment, while meta-learners estimate a binary ATE. Evaluating both against the same continuous ground-truth $\tau(\mathbf{X})$ is scientifically invalid. \\[6pt]
    E3 & 105 & \textcolor{medsev}{Med} &
        Even if the import were fixed, \texttt{DRLearner(models=base)} may not match the constructor---\texttt{econml.dr.DRLearner} typically requires \texttt{model\_regression} and \texttt{model\_propensity} separately. \\[6pt]
    E4 & 126--127 & \textcolor{medsev}{Med} &
        \texttt{confidence\_intervals} catches all exceptions with bare \texttt{except Exception: pass}. Numerical failures are silently swallowed. \\[6pt]
    E5 & 21--23 & \textcolor{lowsev}{Low} &
        \texttt{\_get\_base\_model} silently returns a default model for unknown names. Config typos go unnoticed. \\
    \bottomrule
\end{tabularx}
\end{center}


%% ═════════════════════════════════════════════════════════════════════════
\newpage
\section{Evaluation --- \texttt{src/evaluation.py}}
\label{sec:evaluation}
%% ═════════════════════════════════════════════════════════════════════════

\subsection{Purpose}
Compute accuracy metrics comparing estimated CATE to the ground-truth CATE.

\subsection{Metrics}

\begin{center}
\begin{tabularx}{\textwidth}{l l X}
    \toprule
    \textbf{Metric} & \textbf{Function} & \textbf{What it Measures} \\
    \midrule
    RMSE          & \texttt{cate\_rmse}       & $\sqrt{\frac{1}{n}\sum(\hat{\tau}_i - \tau_i)^2}$ \\[4pt]
    MAE           & \texttt{cate\_mae}        & $\frac{1}{n}\sum|\hat{\tau}_i - \tau_i|$ \\[4pt]
    Bias          & \texttt{cate\_bias}       & $\frac{1}{n}\sum(\hat{\tau}_i - \tau_i)$ \; (positive = overestimation) \\[4pt]
    Spearman $\rho$ & \texttt{rank\_correlation} & Monotonic rank correlation \\
    Kendall $\tau$  & \texttt{rank\_correlation} & Concordance-based rank correlation \\
    CI Coverage   & \texttt{ci\_coverage}     & Fraction of true CATEs inside the confidence interval \\
    GATES         & \texttt{gates\_analysis}  & Binned estimated vs.\ true CATE by quintile \\
    \bottomrule
\end{tabularx}
\end{center}

\subsection{Key Function: \texttt{evaluate\_all()}}

\begin{lstlisting}[caption={Summary builder: runs all metrics for every estimator.}]
def evaluate_all(cate_dict, true_cate, ci_dict=None, n_bins=5):
    rows = []
    for name, est_cate in cate_dict.items():
        row = {"Estimator": name}
        row["RMSE"]  = cate_rmse(est_cate, true_cate)
        row["MAE"]   = cate_mae(est_cate, true_cate)
        row["Bias"]  = cate_bias(est_cate, true_cate)
        rank_corr = rank_correlation(est_cate, true_cate)
        row["Spearman_r"]  = rank_corr["spearman_r"]
        row["Kendall_tau"] = rank_corr["kendall_tau"]
        if ci_dict and name in ci_dict:
            lower, upper = ci_dict[name]
            row["CI_Coverage"] = ci_coverage(true_cate, lower, upper)
        else:
            row["CI_Coverage"] = np.nan
        rows.append(row)
    return pd.DataFrame(rows).set_index("Estimator")
\end{lstlisting}

\subsection{Error Analysis}

\begin{center}
\begin{tabularx}{\textwidth}{c c c X}
    \toprule
    \textbf{ID} & \textbf{Lines} & \textbf{Sev.} & \textbf{Description} \\
    \midrule
    V1 & 90 & \textcolor{highsev}{\textbf{High}} &
        \texttt{iid\_vs\_dependent\_comparison} uses relative import \texttt{from .dgp import ...}, which requires \texttt{src/} to be a proper package (with \texttt{\_\_init\_\_.py}). Will raise \texttt{ImportError} if called. \\[6pt]
    V2 & 87--114 & \textcolor{lowsev}{Low} &
        \texttt{iid\_vs\_dependent\_comparison} is defined but never invoked in \texttt{main.py}. Dead code. \\[6pt]
    V3 & 45 & \textcolor{lowsev}{Low} &
        GATES \texttt{pd.qcut} may fail if an estimator predicts a constant CATE (only one unique value). \\[6pt]
    V4 & 29--31 & \textcolor{lowsev}{Low} &
        \texttt{spearman\_p} and \texttt{kendall\_p} are computed but never included in the results DataFrame. Statistical significance goes unreported. \\
    \bottomrule
\end{tabularx}
\end{center}


%% ═════════════════════════════════════════════════════════════════════════
\newpage
\section{Visualization --- \texttt{src/visualization.py}}
\label{sec:visualization}
%% ═════════════════════════════════════════════════════════════════════════

\subsection{Purpose}
Generate publication-quality plots. All figures are saved to the \texttt{results/} directory.

\subsection{Plot Functions}

\begin{center}
\begin{tabularx}{\textwidth}{l l X}
    \toprule
    \textbf{Function} & \textbf{Output File} & \textbf{Description} \\
    \midrule
    \texttt{plot\_cate\_heatmap}     & \texttt{cate\_heatmap\_true.png} & 2D heatmap of CATE across depth $\times$ volatility bins \\
    \texttt{plot\_cate\_scatter}     & \texttt{scatter\_<est>.png}      & Estimated vs.\ True CATE with 45\textdegree{} reference line \\
    \texttt{plot\_estimator\_comparison} & \texttt{estimator\_comparison\_*.png} & Horizontal bar chart comparing estimators \\
    \texttt{plot\_gates}            & \texttt{gates\_<est>.png}        & Estimated vs.\ True CATE by quintile bin \\
    \texttt{plot\_partial\_dependence} & \texttt{partial\_dependence\_*.png} & Binned avg.\ CATE vs.\ each covariate \\
    \bottomrule
\end{tabularx}
\end{center}

\subsection{Error Analysis}

\begin{center}
\begin{tabularx}{\textwidth}{c c c X}
    \toprule
    \textbf{ID} & \textbf{Lines} & \textbf{Sev.} & \textbf{Description} \\
    \midrule
    P1 & 53 & \textcolor{medsev}{Med} &
        \texttt{np.random.choice} uses the global random state. Scatter plots are non-reproducible across runs. \\[6pt]
    P2 & 31--32 & \textcolor{medsev}{Med} &
        Heatmap bins ``raw'' \texttt{depth} and \texttt{volatility}, but due to bug F1 these columns may already be winsorised. Axis labels are misleading. \\[6pt]
    P3 & 19 & \textcolor{lowsev}{Low} &
        \texttt{figure\_dpi} from config is ignored; style hardcodes \texttt{savefig.dpi: 300}. \\
    \bottomrule
\end{tabularx}
\end{center}


%% ═════════════════════════════════════════════════════════════════════════
\newpage
\section{Pipeline Orchestration --- \texttt{main.py}}
\label{sec:main}
%% ═════════════════════════════════════════════════════════════════════════

\subsection{Pipeline Steps}

\begin{center}
\begin{tabularx}{\textwidth}{c l l X}
    \toprule
    \textbf{Step} & \textbf{Lines} & \textbf{Action} & \textbf{Detail} \\
    \midrule
    Config & 25--26  & Load config         & Read \texttt{params.yaml} \\
    1/6    & 28--31  & Generate data       & 200k synthetic observations, temporal split \\
    2/6    & 33--36  & Engineer features   & Winsorise, derive, standardise \\
    3/6    & 38--47  & Fit estimators      & CausalForest + 4 meta-learners \\
    4/6    & 49--51  & Estimate CATE       & Predict CATE and CIs on test set \\
    5/6    & 53--66  & Evaluate            & All metrics + GATES analysis \\
    6/6    & 68--75  & Visualise \& save   & Generate plots, write CSV \\
    \bottomrule
\end{tabularx}
\end{center}

\subsection{Error Analysis}

\begin{center}
\begin{tabularx}{\textwidth}{c c c X}
    \toprule
    \textbf{ID} & \textbf{Lines} & \textbf{Sev.} & \textbf{Description} \\
    \midrule
    M1 & 39--41 & \textcolor{medsev}{Med} &
        Only 3 of 7 engineered features are used as estimator covariates. Four features (\texttt{log\_spread}, \texttt{inv\_depth}, \texttt{spread\_depth\_ratio}, \texttt{rolling\_vol}) are computed but discarded. \\[6pt]
    M2 & 51 & \textcolor{lowsev}{Low} &
        CI alpha = $1 - 0.90 = 0.10$. This matches econml's convention---\textbf{correct}. \\[6pt]
    M3 & 34--35 & \textcolor{highsev}{\textbf{High}} &
        Consequence of F1: after \texttt{engineer\_features(train\_raw)}, the \texttt{train\_raw} DataFrame is silently mutated. When \texttt{test\_raw} is later passed to plotting functions, it too has been corrupted. \\[6pt]
    M4 & --- & \textcolor{lowsev}{Low} &
        \texttt{src/} has no \texttt{\_\_init\_\_.py}. The \texttt{sys.path.insert} hack is fragile. \\
    \bottomrule
\end{tabularx}
\end{center}


%% ═════════════════════════════════════════════════════════════════════════
\newpage
\section{Cross-Cutting Issues \& Summary}
\label{sec:summary}
%% ═════════════════════════════════════════════════════════════════════════

\subsection{Issue Severity Summary}

\begin{center}
\renewcommand{\arraystretch}{1.3}
\begin{tabular}{lcl}
    \toprule
    \textbf{Severity} & \textbf{Count} & \textbf{Action} \\
    \midrule
    \textcolor{highsev}{\rule{12pt}{12pt}} \; HIGH   & 5 & Fix before trusting results \\
    \textcolor{medsev}{\rule{12pt}{12pt}} \; MEDIUM  & 8 & Fix for correctness and robustness \\
    \textcolor{lowsev}{\rule{12pt}{12pt}} \; LOW     & 7 & Nice to fix, non-critical \\
    \bottomrule
\end{tabular}
\end{center}

\subsection{All High-Severity Issues}

\begin{highbox}
\begin{tabularx}{\textwidth}{c l X}
    \toprule
    \textbf{ID} & \textbf{File} & \textbf{Issue} \\
    \midrule
    E1  & \texttt{estimators.py:102}  & \texttt{DRLearner} import path is wrong $\rightarrow$ \texttt{ImportError} at runtime \\
    E2  & \texttt{estimators.py:40--49} & Comparing continuous vs.\ binary treatment effects against the same ground truth \\
    F1  & \texttt{features.py:59--60} & \texttt{engineer\_features} mutates the input DataFrame in-place \\
    M3  & \texttt{main.py:34--35}     & Raw data used in visualisation is silently corrupted (consequence of F1) \\
    V1  & \texttt{evaluation.py:90}   & Relative import \texttt{from .dgp} will fail without \texttt{\_\_init\_\_.py} \\
    \bottomrule
\end{tabularx}
\end{highbox}

\subsection{All Medium-Severity Issues}

\begin{medbox}
\begin{tabularx}{\textwidth}{c l X}
    \toprule
    \textbf{ID} & \textbf{File} & \textbf{Issue} \\
    \midrule
    F2  & \texttt{features.py:18--23} & Rolling volatility computed from outcome variable (leakage) \\
    F3  & \texttt{features.py + main.py} & 4 engineered features computed but never used \\
    F4  & \texttt{features.py:59--60} & Winsorisation uses test-set quantiles instead of training bounds \\
    D1  & \texttt{dgp.py:79--83} & ``Temporal'' split on non-temporal data \\
    D3  & \texttt{dgp.py (via eval)} & Shallow \texttt{cfg.copy()} leaves nested dicts shared \\
    E3  & \texttt{estimators.py:105} & \texttt{DRLearner(models=base)} may not match constructor API \\
    E4  & \texttt{estimators.py:126--127} & CI exceptions silently swallowed \\
    P1  & \texttt{visualization.py:53} & Non-reproducible scatter sampling \\
    \bottomrule
\end{tabularx}
\end{medbox}

\subsection{Recommended Fix Order}

\begin{enumerate}[leftmargin=2em]
    \item \textbf{Fix DRLearner import (E1)} --- the pipeline will crash without this fix.
    \begin{lstlisting}[numbers=none]
# Change:
from econml.metalearners import DRLearner
# To:
from econml.dr import DRLearner
    \end{lstlisting}

    \item \textbf{Fix DataFrame mutation (F1 / M3)} --- add \texttt{df = df.copy()} at the top of \texttt{engineer\_features}, before winsorisation.
    \begin{lstlisting}[numbers=none]
def engineer_features(df, cfg, fit_stats=None):
    df = df.copy()      # <-- add this line
    feat_cfg = cfg["features"]
    pct = feat_cfg["winsorize_percentile"]
    ...
    \end{lstlisting}

    \item \textbf{Add \texttt{src/\_\_init\_\_.py} (V1 / M4)} --- create an empty file:
    \begin{lstlisting}[language=bash, numbers=none]
touch src/__init__.py
    \end{lstlisting}

    \item \textbf{Address continuous vs.\ binary CATE comparison (E2)} --- either:
    \begin{itemize}
        \item Compute a separate binary ground-truth CATE for the meta-learners, \emph{or}
        \item Use the same treatment type for all estimators.
    \end{itemize}

    \item \textbf{Reuse training winsorisation bounds (F4)} --- store quantile bounds in \texttt{fit\_stats} and reuse for the test set (same pattern as \texttt{standardize}).

    \item \textbf{Remove or use extra features (F3)} --- either add the four unused columns to the estimator covariates, or remove them from \texttt{engineer\_features}.
\end{enumerate}

%% ═════════════════════════════════════════════════════════════════════════

\end{document}
